\chapter{Information theory on the plane}

\textbf{Purpose:} Arrange what the information-theoretic cryptanalysis literature says against what this tree has measured, and expose the one axis that arrangement turns out to run along. \textbf{Scope:} thirty-one hypotheses graded on collision entropy, the Renyi order axis, and the methodological debts the arrangement makes visible.

The short version: thirty-one hypotheses were graded on collision entropy, which is Renyi order two. Order is a continuous axis running from zero to infinity, each point of it answering a different question about the same distribution, and it has never been varied. Order two was not chosen. It is what a collision counter measures, and a collision counter was what was to hand.

\section{What was fetched}

\subsection{Guessing is governed by order one half, not order two}

Arikan's inequality bounds the expected number of guesses for a user with the optimal guessing sequence:
\[
E[G] \le \tfrac{1}{2}\Big[\sum_k \sqrt{p_k}\Big]^2 + \tfrac{1}{2} = \frac{2^{H_{1/2}(X)} + 1}{2}
\]
with the Renyi entropy of order \(\alpha\) written \(H_\alpha(X) = \frac{1}{1-\alpha}\log \sum_x P(x)^\alpha\), and with the property that carries the argument: \(H_\alpha\) is strictly decreasing in \(\alpha\) unless the distribution is uniform on its support.

Massey's earlier bound is in Shannon entropy and is the reason the field talks about Shannon at all here. Arikan's is sharper and names a \emph{different order}: the quantity governing the cost of searching is \(H_{1/2}\), whose defining sum is \(\sum \sqrt{p}\). That sum is dominated by the many small probabilities and not by the few large ones. \textbf{Search cost is a tail statistic.}

\subsection{Collisions are governed by order two, and balance is that}

Bellare and Kohno define the balance of a hash \(h\) with pre-image sizes \(d_i\) as
\[
\mu(h) = \log_r\Big[\frac{d^2}{d_1^2 + \cdots + d_r^2}\Big]
\]
and note immediately that \(1/r^{\mu(h)}\) is the probability that \(h(a) = h(b)\) for \(a, b\) drawn independently from the domain. Their Proposition 5.2 gives \(0 \le \mu(h) \le 1\), zero exactly for a constant function and one exactly for a regular one, and Corollary 6.4 gives the birthday cost as \(Q_h(c) \le 1 + 2.36\sqrt{2c}\, r^{\mu(h)/2}\).

Balance is collision entropy rescaled: \(\mu(h) = H_2 / \log_2 r\). It is this tree's \(H_2\) in a different coordinate, and it answers the collision question, which is not the search question.

\subsection{Distinguishing cost is a divergence}

Baignères, Junod and Vaudenay put the shape of an optimal distinguisher on a hypothesis-testing footing: the optimal statistic is the log likelihood ratio, and the number of samples needed to separate \(P\) from uniform scales as the reciprocal of a divergence between them, with Kullback-Leibler and Chernoff information as the exact exponent and squared Euclidean distance as the usable approximation. The chi-squared method gives the crisp form, advantage over \(n\) samples bounded by \(\sqrt{n\chi^2/2}\).

This is the principled replacement for the sigma counts this work has been reporting. A sigma count says how surprising one statistic is. A divergence says how many samples any distinguisher whatsoever would need, the number deciding whether a structure is worth anything.

\subsection{Iterated functions lose entropy in a known amount}

Flajolet and Odlyzko work through about twenty parameters of a uniformly random map on \(n\) points. The ones that bite here: the image is \((1 - 1/e)n\), leaving \(1/e = 36.79\%\) of the range unreachable in one step; the giant component is \(0.75788n\); rho and tail lengths go as \(\sqrt{\pi n / 8}\).

The holes are not a defect. \textbf{A random function is supposed to have them, and supposed to have exactly \(1/e\) of them.}

\section{The arrangement: order is the axis}

Every result above is a statement about the same object, the output distribution, read at a different Renyi order.

\begin{center}
\begin{tabular}{@{}llll@{}}
\toprule
order & reads & what it answers & who \\
\midrule
\(\alpha = 0\) & size of the support & how many values never occur, the holes & Flajolet-Odlyzko \\
\(\alpha = 1/2\) & \(\sum\sqrt{p}\) & cost of guessing, the search & Arikan \\
\(\alpha = 1\) & Shannon & information content & Massey \\
\(\alpha = 2\) & collision probability & cost of colliding, balance & Bellare-Kohno \\
\(\alpha = \infty\) & largest single probability & the biggest bump & min-entropy \\
\bottomrule
\end{tabular}
\end{center}

Holes at one end, bumps at the other, guessing and collisions at two interior points. The plane is one plane; the order is the angle it is viewed from. This tree has looked at it from one angle.

\subsection{What a flat plane predicts along that axis}

Write the bin counts as \(c_i = m(1 + e_i)\) with \(m = d/r\) the exact mean. For a random function the \(e_i\) are small and the deficit \(\log_2 r - H_\alpha\) expands to
\[
\mathrm{deficit}(\alpha) = \frac{\alpha\,\mathrm{var}(e)}{2\ln 2}, \qquad \mathrm{var}(e) = \frac{r-1}{d} \text{ exactly}
\]

The deficit profile of a random function is \textbf{a straight line through the origin in \(\alpha\)}, which gives a test with no free parameter at all:
\[
\mathrm{deficit}(2) / \mathrm{deficit}(1/2) = 4 \quad\text{exactly}
\]

Every scale cancels out of a ratio between two orders: domain size, bin count, window position, how many samples were taken. Curvature in that line, or a ratio away from four, is structure that cannot be explained away by the size of the experiment.

Where the mean count falls to one the expansion stops applying and the counts go Poisson, whose moments about a unit mean are the Bell numbers \(1, 2, 5, 15\). Both regimes are real here: a 32-bit window read over a \(2^{32}\) domain is the Poisson case, and it is the only place a bin can be empty. That is where the holes live.

\section{What this says about work already done}

\subsection{Every advantage quoted from an \texorpdfstring{\(H_2\)}{H2} measurement overstates the search advantage}

\(H_\alpha\) is decreasing in \(\alpha\). Therefore \(H_{1/2} \ge H_2\) and \(\mathrm{deficit}(1/2) \le \mathrm{deficit}(2)\), the first being about a quarter of the second for anything near flat. Every figure of the form \emph{collision entropy is \(X\) bits below uniform, therefore there are \(X\) bits of advantage} is an answer to the collision question read as an answer to the search question, and it is the more generous of the two.

The concrete case: H22's rotational coherence measured \(H_2 = 16.503\) against a uniform \(32\), a deficit of \(15.497\) bits, with a permuted control at \(32.069\). That deficit is real and the control is sound. What it does not say is that \(15.497\) bits of search advantage exist, because search is governed by order one half. The honest statement is \(\mathrm{deficit}(1/2) \le 15.497\), and the actual value needs the same histogram re-read at order one half. \textbf{That re-read is cheap and has not been done.}

This is not a small correction of presentation. It is the difference between the quantity measured and the quantity the project is about.

\subsection{Bellare and Kohno's own SHA-1 numbers have an unsubtracted null}

Their full version computes the balance of SHA-1 restricted to a \(2^{32}\) domain at byte-aligned output windows and reports the values are high. Subtracting the random-function null instead of eyeballing the values:

\begin{center}
\begin{tabular}{@{}lrrr@{}}
\toprule
window & their measured deficit & random function predicts & ratio \\
\midrule
8 bits & \(1.0462 \times 10^{-8}\) & \(1.0707 \times 10^{-8}\) & \(0.9771\) \\
16 bits & \(1.3770 \times 10^{-6}\) & \(1.3758 \times 10^{-6}\) & \(1.0009\) \\
24 bits & \(2.34378 \times 10^{-4}\) & \(2.34356 \times 10^{-4}\) & \(1.0001\) \\
\bottomrule
\end{tabular}
\end{center}

Their own resolution at width 8 is \(\sqrt{2/255} = 8.8\%\), and \(0.977\) sits well inside it. \textbf{All of SHA-1's measured imbalance is the imbalance any function drawn at random would show, to four significant figures.} The paper says the balance is high. The sharper statement available from the same table is that the deficit is exactly generic, and that is a much stronger claim about SHA-1 than \emph{high}.

It is also the methodological point this tree keeps relearning: the interesting quantity is never the measurement, it is the measurement divided by what a random function would have given.

\subsection{Sigma counts should be divergences}

Comparing a max-statistic against a per-bin uniform and getting \(775{,}000\times\) where the sigma-delta null gives \(184.5\times\) is the same error the distinguishing literature solved in 2004. A divergence is the right currency because it converts directly into the number of samples any distinguisher would need, the only thing deciding whether a structure is exploitable. This is not yet done and it is the largest unresolved methodological debt.

\section{The compiler audit}

A compiler folds, contracts, reassociates and hoists, and a folded measurement is indistinguishable by eye from a computed one. Results depending on the optimizer are not results.

Two arms, both mechanical. \textbf{Differential}: build every bench three ways, \texttt{-O0}, \texttt{-O2}, and \texttt{-O2 -ffp-contract=off}, and run all three. A bench whose three outputs are byte-identical computed its numbers instead of inheriting them, and since GCC contracts multiply-add into FMA by default, the third arm separates that from folding proper. \textbf{Constant pool}: read every double in a binary's read-only data and report which of the numbers the bench printed were already sitting there before it ran.

Neither arm proves a folded constant was folded \emph{correctly}, a third failure mode, and for the Renyi bench that is closed separately by reimplementing every printed prediction in another language against another libm.

\section{What the posits say about all of this}

Four posits are proved and not asserted, each on constructed domains where the answer is known by design instead of inferred from the case that suggested it. Three bear directly here and one reframes this entire chapter.

\subsection{Histogram quantities are free, and everything above is a histogram quantity}

Histogram quantities are permutation invariant. They describe the maximum entropy case and are free. The arrangement is what is left after that.

\textbf{Every number the Renyi bench produces is a histogram quantity.} All seven Renyi orders, the Bellare-Kohno balance, the count of holes, the Arikan guessing factor: each is a function of the multiset of bin counts and of nothing besides. Shuffling which value carries which count leaves every one of them exactly unchanged.

So the headline result, that SHA256d's output histogram over a completely enumerated \(2^{32}\) nonce domain is that of a random function to six significant figures, is entailed instead of discovered. It is the maximum-entropy case and it was always going to be the answer. That is worth having as a measured fact instead of an assumption, and it is not a finding about the arrangement, because no permutation-invariant quantity can be.

\subsection{The null must delete the property being asked about}

Proved on four constructed clouds, lattice or scatter crossed with patterned or random values, where a value-permuting null can only reveal order in the values and a position-moving null can only reveal order in the geometry. The predicted hits therefore fall on a diagonal, and anything off it refutes the posit.

The permutation null over arrangements of a fixed multiset satisfies this exactly for the question asked here: it preserves every count and destroys every position. It deletes the arrangement and keeps the counts. Drawing uniformly from the arrangements of a fixed multiset is the least committal distribution consistent with the observed histogram. It asserts nothing beyond the quantity already measured. Every other background in this work is a model and can be false. This one is the data with one property deleted.

\subsection{A quantity that fails to conserve indicates the instrument}

State what a measure must be invariant to and check each, instead of inferring the rule from the case that suggested it. Written out for the deficit measured here:

\begin{center}
\begin{tabular}{@{}lp{0.42\linewidth}@{}}
\toprule
transformation & required \\
\midrule
relabel the window values by any fixed bijection & exactly invariant, the histogram is unchanged \\
reverse or permute the order of nonce enumeration & exactly invariant, counts do not see order \\
split the domain and recombine the counts & exactly invariant \\
read a different window position & invariant within the stated spread \\
\textbf{run on a different block header} & invariant within the stated spread, \textbf{never checked} \\
\bottomrule
\end{tabular}
\end{center}

The last row is the gap. Every measurement in this tree is on one header. One header is one corpus, and the homogeneity check the posit demands has not been run. It is the same shape as the one-seed failure mode applied to the header instead of to the seed.

The posit's other half is sharper still: the spread over seeds is the floor below which no difference between two corpora means anything, and that floor has been used all through this work and never measured. Here the ratios \(1.00967\) for SHA and \(0.98900\) for splitmix differ by about \(2\%\) at the byte windows, against an analytic per-family floor of \(1.57\%\). That comparison is sound only because the floor is derivable exactly for a multinomial. Where a floor is not derivable it has to be measured by reseeding, and that has not been done anywhere in this tree.

\subsection{The symbol width has to match the scale of the structure}

A detector is not told where the units begin, and a slice of the right width at the wrong offset splits every unit across two symbols. Both the width and the phase have to be swept.

Every window in the Renyi bench is byte-aligned: 8 bits at 32 positions, 16 at 16, 32 at 2. The phase sweep has since been run at eight phases, and the sixteen-bit family reads a spread ratio of \(0.997\) against its control. No alignment dependence.

\subsection{The content has to be printed, not only the statistic}

Nine problems in the companion work were found by reading output and none by a statistic going out of range. This tree prints statistics almost exclusively. The exception is recent and earned its keep immediately: extending the leading-zero table past 44 bits showed that the single deep nonce is the block's own solved nonce, rediscovered by the enumeration. That is the strongest external check available here, and it had been sitting one row below the cut.

\subsection{The cost term is the finding}

Widening the search for structure in the hole arrangement, one step at a time:

\begin{center}
\begin{tabular}{@{}p{0.26\linewidth}rp{0.2\linewidth}p{0.24\linewidth}@{}}
\toprule
what was checked & directions & what it cost & what it found \\
\midrule
single-bit flips of the count field & 32 & one pass over 4 GB, seconds & nothing, max \(2.28\sigma\) \\
every mask inside any two bytes & 393,210 & seconds & nothing, max \(5.03\sigma\) against a control at \(4.82\) \\
\textbf{every mask there is} & \(4{,}294{,}967{,}295\) & 16 GB, \(1.4\times10^{11}\) butterflies, minutes & nothing, max \(6.79\sigma\) against an expected \(6.24\) \\
\bottomrule
\end{tabular}
\end{center}

Each step multiplies the cost by orders of magnitude and returns the same answer. The free half was free: the output histogram matches a random function's to six significant figures and was always going to. The other half is made of checks, and the number of checks scales with the space being checked. Testing every linear direction over \(2^{32}\) values costs \(n\log n\); testing every quadratic one costs more; the set of structures that could exist is larger than the space itself by an amount no budget closes.

\textbf{This is the same asymmetry that makes mining work, arriving from the other side.} Knowing the function completely removes none of the checks. That is worth stating plainly because it answers \emph{we know the hash, we know what they did to it}: knowing the transform does not reduce the count of things that have to be verified, and the count is the cost.

\subsection{The knee, and where it sits}

The cost does not merely grow. It floors, with a knee, and the knee's position is derivable instead of a matter of opinion.

Checking \(N\) directions means taking a maximum over \(N\) statistics. The threshold a real signal has to clear is the expected largest of \(N\) nulls. Reach goes as \(R(N) = \sqrt{2\ln N}\) minus a slowly varying correction, and cost as \(C(N) \sim N\log N\) in time and \(N\) in memory. Inverting and substituting gives \(N = e^{R^2/2}\) and \(C \sim e^{R^2/2}(R^2/2)\).

\textbf{The cost is exponential in the square of the reach.}

\begin{center}
\begin{tabular}{@{}lrl@{}}
\toprule
reach wanted & directions needed & relative to what was run \\
\midrule
\(4.57\sigma\) & 393,210 & 1/10,900 \\
\textbf{\(6.24\sigma\)} & \(4.29\times10^{9}\) & \textbf{1, and it cost 16 GB and minutes} \\
\(8\sigma\) & \(\sim 4.3\times10^{15}\) & \(\sim 10^{6}\times\), about 17 petabytes \\
\(10\sigma\) & \(\sim 7.7\times10^{23}\) & not on this planet \\
\bottomrule
\end{tabular}
\end{center}

Going from every two-byte mask to every mask was \(10{,}900\) times the work and bought \(1.67\sigma\). The next \(1.67\sigma\) costs about a million times more than that, and the one after costs ten billion times more again. The complete-direction approach floors at about \(7\sigma\) for any budget that exists.

\textbf{The knee is already behind us, and the way past it is not more directions.} More directions is the move whose price is exponential in the square of what it buys. Getting further needs a statistic where a signal, if there is one, is concentrated in a few places known in advance instead of spread across billions of equally plausible ones, and that is a statement about having a hypothesis and not about having a bigger machine.

\section{Still unknown}

\textbf{A second header.} One block's header is one corpus, and nothing here has been asked to conserve across a different one. Cheap now that the enumeration is on the device.

\textbf{A reduced-width analog.} An 8-bit-word SHA-256 has a state small enough to enumerate completely. That is the axis where speed still converges the representation, because it replaces every sampled result with an exact one instead of buying another \(\sqrt{2\ln N}\) of reach against an exponential price. It is the only open item that is a better object instead of a bigger search.

\textbf{The order one half re-read} of every histogram already collected. Cheap, and it changes what the existing advantage figures mean.

\textbf{Divergences in place of sigma counts}, throughout.

\textbf{The functional graph proper.} Image size, rho and tail lengths, component structure of the compression function as a map on a restricted domain, against Flajolet-Odlyzko's constants.

\textbf{Whether the deficit line stays straight at reduced rounds.} Round eight is where closeness dies; the order axis has never been swept at a partial round, and a curvature that appears and then flattens would locate the mixing in a coordinate nothing here uses.

\section{Sources}

\begin{enumerate}
  \item E. Arikan, \emph{An Inequality on Guessing and Its Application to Sequential Decoding}, IEEE Trans. Inf. Theory 42(1):99--105, 1996. Restated with the constants in IACR ePrint 2018/455.
  \item J. Massey, \emph{Guessing and Entropy}, ISIT 1994.
  \item M. Bellare and T. Kohno, \emph{Hash Function Balance and its Impact on Birthday Attacks}, EUROCRYPT 2004. IACR ePrint 2003/065.
  \item T. Baignères and S. Vaudenay, \emph{The Complexity of Distinguishing Distributions}, ICITS 2008.
  \item W. Dai, V. T. Hoang and S. Tessaro, \emph{Information-theoretic Indistinguishability via the Chi-squared Method}. IACR ePrint 2017/537.
  \item P. Flajolet and A. Odlyzko, \emph{Random Mapping Statistics}, EUROCRYPT 1989.
  \item \emph{Iteration Entropy}, arXiv 1712.01407.
  \item D. Malone and W. Sullivan, \emph{Guesswork is not a Substitute for Entropy}, ITT 2005.
\end{enumerate}
